% Generated draft: requires mathematical and source review.
\documentclass[11pt]{article}
\usepackage[margin=25mm]{geometry}
\usepackage{fontspec}
\setmainfont{Latin Modern Roman}
\usepackage{amsmath,amssymb,mathtools,unicode-math}
\setmathfont{Latin Modern Math}
\usepackage{enumitem,needspace,xurl,hyperref}
\setlength{\parindent}{0pt}
\setlength{\parskip}{5pt}
\setlength{\emergencystretch}{3em}
\newcommand{\N}{\mathbb{N}}
\newcommand{\Z}{\mathbb{Z}}
\newcommand{\Q}{\mathbb{Q}}
\newcommand{\R}{\mathbb{R}}
\newcommand{\C}{\mathbb{C}}
\newcommand{\F}{\mathbb{F}}
\newcommand{\A}{\mathbb{A}}
\newcommand{\PP}{\mathbb{P}}
\newcommand{\E}{\mathbb{E}}
\newcommand{\eps}{\varepsilon}
\newcommand{\dd}{\,\mathrm{d}}
\begin{document}
\section*{UnsolvedMath ID: 30002914}
\textbf{Source problem code:} catalogue entry
\subsection{Definitions and mathematical statement}
For a smooth Riemannian manifold \((\Sigma,h)\), its Riemannian cone away from the vertex is
\[
C(\Sigma)\setminus\{o\}=(0,\infty)\times\Sigma,
\qquad g_C=\dd r^2+r^2h.
\]
The vertex is supplied by metric completion when appropriate. This notation describes the smooth-link case; the source does not specify whether only such cones are intended, nor does it specify dimension, compactness of the link, boundary conditions, or permissible singularities.

A gradient Ricci expander is a smooth Riemannian manifold \((M,g)\) with a smooth function \(f\colon M\to\R\) satisfying
\[
\operatorname{Ric}_g+\nabla_g^2f=\lambda g,
\qquad \lambda<0.
\]
Here \(\operatorname{Ric}_g\) is the Ricci tensor and \(\nabla_g^2f\) is the Hessian. This uses the potential convention in which expansion corresponds to a negative soliton constant.

The supplied problem asks:
\begin{enumerate}
\item Which metric cones admit a smoothing by a gradient Ricci expander asymptotic to the cone?
\item Given a conical gradient Ricci expander, can it be deformed by an implicit-function construction, including under suitable variations of its asymptotic cone?
\item Are the resulting deformations globally unique?
\end{enumerate}
The interpretation of smoothing through an asymptotic cone is the geometric framework suggested by the wording; the source does not define it. In particular, it supplies no completeness requirement, convergence topology or decay rate, curvature hypotheses, deformation spaces, or equivalence relation for uniqueness. It also does not specify whether global uniqueness concerns all expanders with a prescribed cone or continuation of a local deformation family. Thus these are qualified research questions, and the supplied evidence does not determine a single fully quantified existence, deformation, or uniqueness assertion.
\subsection{Short English statement}
Determine which metric cones admit gradient Ricci expander smoothings, whether conical expanders admit implicit deformations, and whether those deformations are globally unique. The source leaves the admissible cones, asymptotic conditions, deformation spaces, and meaning of global uniqueness unspecified; these require mathematical review.
\subsection{Sources}
\begin{itemize}
\item Local UnsolvedMath registry, supplied complete source record: permanent ID \(30002914\), title \emph{Smoothing and Deforming Metric Cones by Ricci Expanders}, source problem code \texttt{catalogue entry}, in the Oberwolfach Reports Open Problems collection. No individual proposer is supplied; the proposed year is recorded as 2015.
\item Supplied report citation: \emph{Differentialgeometrie im Großen (2016)}, Oberwolfach Reports, pp.~1759--1807, DOI \texttt{10.4171/owr/2015/31}. Original Oberwolfach report: \url{https://doi.org/10.4171/owr/2015/31}. The supplied title year and proposed year differ; their relationship has not been verified.
\item Supplied literature reference labeled ``Deruelle primary paper'': \url{https://doi.org/10.1093/imrn/rnw110}. No paper title or precise theorem attribution is supplied.
\item Supplied literature reference labeled ``Conlon and Deruelle primary paper'': \url{https://arxiv.org/abs/1607.03546}. No paper title or precise theorem attribution is supplied.
\item Literature assessment dated exactly \texttt{2026-08-21}: the supplied curation labels the problem partially solved and reports implicit deformation and unique-continuation results for important conical Ricci expander classes, while global uniqueness and unrestricted deformation remain open. This is supplied dated curation derived from the Oberwolfach Reports record, not an independently verified literature assessment. Its time-sensitive status and cited claims require independent verification.
\end{itemize}
\end{document}
